%=====================================================================
\section{Specification of the components not yet implemented}\label{sec:spec:top}
\label{sec:spec:pending}

\begin{quote}
\textbf{Read this before any equation below.} \emph{Nothing in this section is implemented in the
current model version.} The four components specified here --- adiabatic humidifier, indirect
evaporative cooler, single-effect LiBr--water absorption chiller, wet cooling tower --- exist in the
master architecture and in the mode table, but no line of code evaluates them. Everywhere else in
this report the equations were read off the notebook and the numbers harvested from it, so the text
cannot disagree with the model. Here the direction is reversed: this is the \emph{design document},
written so that implementing it is transcription rather than invention. Consequently every numerical
value quoted below is a \emph{typical literature value}, offered to make the equations concrete and
to give the implementer a sanity range. It is not a code constant, it has been harvested from
nothing, and it must be checked against a real machine or a real correlation before it sizes
hardware. Each value is flagged where it appears.
\end{quote}

Three of the four are cheap: one algebraic block each. The fourth --- the chiller, which cannot be
implemented without the tower and the chilled-water loop --- is the expensive one, and it is also the
one that decides whether the well can serve the cooling corner at all.
Section~\ref{sec:spec:order} closes with the implementation order and with the one architectural
defect this exercise has exposed. Throughout, air states are carried as $(T,w)$ pairs and enthalpies
use the notebook's own correlation, $h = 1.006\,T + w(2501 + 1.86\,T)$ in
\si{\kilo\joule\per\kilogram} of dry air, so the new blocks compose with the existing ones without a
units seam. Each block returns a dictionary of outlet states and duty, matching the pattern of
\texttt{regen\_coil} and \texttt{desiccant\_wheel}.

\paragraph{A prerequisite that is not a component.}
Section~\ref{sec:psy:wetbulb} recorded that the notebook contains \emph{no} wet-bulb routine and no
dew-point routine, because nothing in the present solution path needs one. That ceases to be true
the moment evaporative equipment enters. The implementation must therefore first add two scalar
solvers,
\begin{equation}
h\!\left(T_{\mathrm{wb}},\, w_{s}(T_{\mathrm{wb}})\right) - h(T,w)
 = \left(w_{s}(T_{\mathrm{wb}}) - w\right) c_{p,\mathrm{w}}\,T_{\mathrm{wb}} ,
\qquad
\psat(T_{\mathrm{dp}}) = \frac{w\,\patm}{0.622 + w} ,
\label{eq:spec:wbdp}
\end{equation}
with $w_{s}(T) = 0.622\,\psat(T)/(\patm-\psat(T))$ and
$c_{p,\mathrm{w}}=\SI{4.186}{\kilo\joule\per\kilogram\per\kelvin}$. The left is the
adiabatic-saturation (thermodynamic wet-bulb) definition; the right inverts the Magnus correlation
\eqref{eq:psy:psat} in closed form. Both are monotone and safe for bisection. The makeup term on the
right of the first shifts $T_{\mathrm{wb}}$ by at most \SIrange{0.1}{0.2}{\kelvin} at the humidities
of interest, which is why ``constant wet bulb'' and ``constant enthalpy'' are used interchangeably
below.

%---------------------------------------------------------------------
\subsection{Adiabatic humidifier (unlocks M4)}
\label{sec:spec:hum}

\emph{Status: specification only. Not implemented in the current model version.}

\subsubsection{Hypothesis and equations}
Air passes over a wetted surface --- rigid cellulose media, or an atomising spray --- fed with
recirculated water that has equilibrated at the air's own wet bulb. No heat crosses the casing and no
work is done, so the energy to evaporate the water can come from one place only: the sensible heat of
the air itself. The process is adiabatic saturation, and on the chart it walks the air \emph{up the
constant wet-bulb line} toward saturation.

With inlet $(T_{\mathrm{hi}}, w_{\mathrm{hi}})$ and its wet bulb $T_{\mathrm{hi}}^{\mathrm{wb}}$ from
\eqref{eq:spec:wbdp}, the device is characterised by one number, the \emph{saturation
effectiveness} --- the fraction of the available wet-bulb depression the media actually delivers:
\begin{equation}
\boxed{\;
\varepsilon_{\mathrm{sat}} \equiv
\frac{T_{\mathrm{hi}} - T_{\mathrm{ho}}}{T_{\mathrm{hi}} - T_{\mathrm{hi}}^{\mathrm{wb}}}
\;}
\qquad\Longrightarrow\qquad
T_{\mathrm{ho}} = T_{\mathrm{hi}} - \varepsilon_{\mathrm{sat}}
 \left(T_{\mathrm{hi}} - T_{\mathrm{hi}}^{\mathrm{wb}}\right) .
\label{eq:spec:epssat}
\end{equation}
The outlet humidity ratio is \emph{not} a second free parameter: it is fixed by the adiabatic
constraint $h_{\mathrm{ho}} = h_{\mathrm{hi}}$, inverted with the notebook's enthalpy correlation,
\begin{equation}
w_{\mathrm{ho}} = \frac{h_{\mathrm{hi}} - 1.006\,T_{\mathrm{ho}}}{2501 + 1.86\,T_{\mathrm{ho}}} ,
\qquad
h_{\mathrm{hi}} = 1.006\,T_{\mathrm{hi}} + w_{\mathrm{hi}}(2501 + 1.86\,T_{\mathrm{hi}}) .
\label{eq:spec:humw}
\end{equation}
The water consumed is the moisture actually added,
\begin{equation}
\md_{\mathrm{w,hum}} = \mdp\left(w_{\mathrm{ho}} - w_{\mathrm{hi}}\right) ,
\label{eq:spec:humwater}
\end{equation}
plus a sump bleed of \SIrange{10}{25}{\percent} of \eqref{eq:spec:humwater} for recirculating media
(a \emph{typical literature figure}, not a code constant), and zero for a once-through atomiser. The
block returns $\{T_{\mathrm{ho}}, w_{\mathrm{ho}}, \md_{\mathrm{w,hum}},
\varepsilon_{\mathrm{sat}}\}$ and, being adiabatic, no duty. Assert
$w_{\mathrm{ho}} \le w_{s}(T_{\mathrm{ho}})$ on exit: with $\varepsilon_{\mathrm{sat}}<1$ it cannot
fail, and if it ever does the enthalpy inversion was called with an inconsistent state.

\paragraph{Typical values.} Rigid cellulose \SIrange{200}{300}{\milli\metre} deep at
\SI{2.5}{\metre\per\second} face velocity gives $\varepsilon_{\mathrm{sat}} = \numrange{0.85}{0.95}$;
atomising and ultrasonic units run \numrange{0.70}{0.90}. Use $0.90$ as a placeholder.
\emph{Literature ranges quoted for concreteness; not harvested from the model, because the model does
not contain this component.}

\subsubsection{What it means physically, and why M4 needs no well heat}
Equation~\eqref{eq:spec:humw} says the enthalpy the air gains as vapour it loses as temperature. As a
duty balance across the device,
\begin{equation}
\qs^{\mathrm{del}} = \mdp\,c_{p,\mathrm{a}}\left(T_{\mathrm{hi}}-T_{\mathrm{ho}}\right),
\qquad
\ql^{\mathrm{del}} = -\,\mdp\,\hfg\left(w_{\mathrm{ho}}-w_{\mathrm{hi}}\right),
\qquad
\qs^{\mathrm{del}} + \ql^{\mathrm{del}} \approx 0 .
\label{eq:spec:humduty}
\end{equation}
This is the whole point of mode M4. \emph{Adiabatic saturation is direct evaporative cooling}: the
humidifier is not a component that humidifies while something else cools, it is one device
discharging both requirements of the cooling-plus-humidification quadrant simultaneously, and it does
so by moving heat \emph{within} the air stream rather than importing any. No geothermal water is
drawn, no wheel turns, no chiller runs. That is why the mode table records ``water required: none''
for M4 --- the one mode in the architecture that is free of the reservoir entirely, and therefore the
one mode that survives unchanged for the whole plant life.

\paragraph{The constraint that comes with the gift.} Equation~\eqref{eq:spec:humduty} is also a
restriction. Because the process is pinned to a constant-enthalpy line, the ratio of sensible to
latent delivery is not a design variable: it is whatever the psychrometrics of that line happen to
be. The device has \emph{one} degree of freedom --- $\varepsilon_{\mathrm{sat}}$, or a bypass damper
--- and one degree of freedom can satisfy $\qs$ or $\ql$, not both. Modulation slides the outlet
along the same ray from the inlet state; it does not rotate that ray. M4 therefore serves the
sub-region of its quadrant whose load ratio lies near that ray, and loads far off it fall through to
M5, where a chilled-water coil supplies the missing sensible capacity and the humidifier trims the
latent. That handoff is where the ordering defect of Section~\ref{sec:spec:order} lives.

%---------------------------------------------------------------------
\subsection{Indirect evaporative cooler and the M-cycle (unlocks M2)}
\label{sec:spec:iec}

\emph{Status: specification only. Not implemented in the current model version.}

\subsubsection{Hypothesis}
Two streams are separated by an impermeable plate. The \emph{primary} (product) stream flows on the
dry side and is cooled sensibly; the \emph{secondary} (working) stream flows on the wet side,
evaporates water into itself, and is vented. Because the plate is impermeable the product gains no
moisture,
\begin{equation}
w_{\mathrm{1o}} = w_{\mathrm{1i}} ,
\label{eq:spec:iecw}
\end{equation}
which is the property that makes this, and not a direct evaporative cooler, the correct partner for
the desiccant wheel. The wet side is taken as saturated at exit.

\subsubsection{Effectiveness, in two flavours}
In a \emph{conventional} IEC the working air enters the wet channel at the inlet state, so its wet
bulb is the inlet wet bulb and that is the coldest sink available:
\begin{equation}
\boxed{\;
\varepsilon_{\mathrm{wb}} \equiv
\frac{T_{\mathrm{1i}}-T_{\mathrm{1o}}}{T_{\mathrm{1i}}-T_{\mathrm{1i}}^{\mathrm{wb}}}\;}
\qquad
T_{\mathrm{1o}} = T_{\mathrm{1i}} - \varepsilon_{\mathrm{wb}}
 \left(T_{\mathrm{1i}}-T_{\mathrm{1i}}^{\mathrm{wb}}\right) .
\label{eq:spec:iecwb}
\end{equation}
In the \emph{M-cycle} (Maisotsenko) regenerative arrangement the working air is first pre-cooled
sensibly along a dry channel and only then turned through perforations into the adjacent wet channel.
Each successive element of working air therefore enters the wet side already colder, with a depressed
wet bulb, and the asymptotic sink is not the inlet wet bulb but the inlet \emph{dew point}. The
effectiveness must then be referred to the dew-point depression:
\begin{equation}
\boxed{\;
\varepsilon_{\mathrm{dp}} \equiv
\frac{T_{\mathrm{1i}}-T_{\mathrm{1o}}}{T_{\mathrm{1i}}-T_{\mathrm{1i}}^{\mathrm{dp}}}\;}
\qquad
T_{\mathrm{1o}} = T_{\mathrm{1i}} - \varepsilon_{\mathrm{dp}}
 \left(T_{\mathrm{1i}}-T_{\mathrm{1i}}^{\mathrm{dp}}\right) .
\label{eq:spec:iecdp}
\end{equation}
The two are not interchangeable, and confusing them is the classic error in this literature: since
$T^{\mathrm{dp}} < T^{\mathrm{wb}} < T$, a machine with $\varepsilon_{\mathrm{dp}}=0.6$ can exhibit an
\emph{apparent} $\varepsilon_{\mathrm{wb}}$ above unity. That violates nothing; it says only that the
wet bulb was never the correct reference. The block should accept both and record which was used.

\subsubsection{Outlets, duty and water}
Let $\md_{\mathrm{in}}$ be the air delivered, split by the \emph{working-air fraction} $\phi$:
\begin{equation}
\md_{\mathrm{sec}} = \phi\,\md_{\mathrm{in}} ,
\qquad
\md_{\mathrm{pri}} = (1-\phi)\,\md_{\mathrm{in}} .
\label{eq:spec:iecphi}
\end{equation}
The primary outlet is \eqref{eq:spec:iecw} with \eqref{eq:spec:iecwb} or \eqref{eq:spec:iecdp}, and
the duty removed from the product stream is
\begin{equation}
Q_{\mathrm{IEC}} = \md_{\mathrm{pri}}\,c_{p,\mathrm{ma}}\left(T_{\mathrm{1i}}-T_{\mathrm{1o}}\right),
\qquad c_{p,\mathrm{ma}} = 1.006 + 1.86\,w_{\mathrm{1i}} .
\label{eq:spec:iecQ}
\end{equation}
(The existing coil block uses the dry-air value \texttt{CP\_A} in the same place; at
$w\approx\SI{10}{\gram\per\kilogram}$ the difference is under \SI{2}{\percent}. Whichever is chosen,
choose it once and use it in both.) The secondary stream carries that heat away and leaves saturated;
its state follows from an energy balance closed against the saturation line --- one monotone
root-find in $T_{\mathrm{2o}}$:
\begin{equation}
\md_{\mathrm{sec}}\left[h\!\left(T_{\mathrm{2o}},w_{s}(T_{\mathrm{2o}})\right)-h_{\mathrm{2i}}\right]
= Q_{\mathrm{IEC}} + \md_{\mathrm{w,IEC}}\,c_{p,\mathrm{w}}\,T_{\mathrm{sump}} ,
\qquad
w_{\mathrm{2o}} = w_{s}(T_{\mathrm{2o}}) ,
\label{eq:spec:iecsec}
\end{equation}
with water consumption
\begin{equation}
\md_{\mathrm{w,IEC}} = \md_{\mathrm{sec}}\left(w_{\mathrm{2o}}-w_{\mathrm{2i}}\right) .
\label{eq:spec:iecwater}
\end{equation}
Dropping the makeup-enthalpy term costs a few tenths of a kelvin in $T_{\mathrm{2o}}$ and makes
\eqref{eq:spec:iecsec} explicit; retaining it costs one iteration and is cleaner. The block returns
both outlet states, the duty $Q_{\mathrm{IEC}}$, the water rate $\md_{\mathrm{w,IEC}}$ and --- this
is the one easy to forget --- the \emph{reduced} product flow $\md_{\mathrm{pri}}$, which the
downstream mass balance must honour.

\paragraph{Typical values.} Conventional IEC $\varepsilon_{\mathrm{wb}}=\numrange{0.55}{0.75}$;
M-cycle $\varepsilon_{\mathrm{dp}}=\numrange{0.55}{0.75}$, i.e.\ an apparent wet-bulb effectiveness of
roughly \numrange{1.0}{1.2}; working-air fraction $\phi=\numrange{0.30}{0.40}$. \emph{Literature
ranges, not code constants.}

\subsubsection{Why the M-cycle is the indicated partner for the wheel}
The wheel leaves the process air \emph{hot and dry} --- the central consequence recorded in
Section~\ref{sec:dw:latent}: latent load has been converted into sensible load. Three things follow,
and together they select the M-cycle.
\begin{enumerate}[nosep]
\item \textbf{The moisture margin must be protected.} After the wheel $w$ already sits below the
supply target, and regeneration heat was spent to put it there. A direct evaporative cooler hands
that margin straight back as vapour; the IEC cools at constant $w$ by construction,
Eq.~\eqref{eq:spec:iecw}, and spends none of it.
\item \textbf{Dry air has a deep dew point.} The depression $T_{\mathrm{1i}}-T_{\mathrm{1i}}^{
\mathrm{dp}}$ that \eqref{eq:spec:iecdp} multiplies grows rapidly as $w$ falls. The wheel does not
merely leave a load for the M-cycle to remove; it \emph{enlarges the M-cycle's driving potential} in
the act of creating that load. The two are complementary, not merely sequential.
\item \textbf{Sub-wet-bulb is what makes the target reachable.} A conventional IEC bottoms out at
$T^{\mathrm{wb}}$, which for post-wheel air may still sit above the supply setpoint. The M-cycle's
floor is $T^{\mathrm{dp}}$, several kelvin lower, and that margin is what turns M2 from marginal into
comfortable.
\end{enumerate}

\paragraph{The cost of $\phi$.} The working-air fraction is the design lever and it cuts both ways.
Raising $\phi$ adds wet-channel surface and enthalpy sink per unit of product, so both effectivenesses
rise and $T_{\mathrm{1o}}$ falls --- but the product flow $(1-\phi)\md_{\mathrm{in}}$ falls with it,
and the delivered capacity $Q_{\mathrm{IEC}}\propto(1-\phi)\,\varepsilon(\phi)$ is a rising factor
times a falling one. It has an interior maximum, typically near $\phi\approx\numrange{0.3}{0.4}$,
so the design question is not ``how cold can the product get'' but ``at what $\phi$ does the
product \emph{duty} peak''. This is the coil's temperature-versus-throughput tension
(Section~\ref{sec:hyd:coil}) in a different currency.

%---------------------------------------------------------------------
\subsection{Single-effect LiBr--water absorption chiller (unlocks M3 and M5)}
\label{sec:spec:abs}

\emph{Status: specification only. Not implemented in the current model version, and it cannot be
implemented alone --- the chilled-water loop and the tower of Section~\ref{sec:spec:tower} belong to
the same step.}

\subsubsection{Hypothesis: a characteristic equation, not a cycle model}
A full cycle model carries solution concentrations, a solution heat exchanger and crystallisation
margins, and needs property data for aqueous lithium bromide. None of that is warranted here, because
the plant-level question is only how much cooling the machine delivers from a given hot-water
temperature at a given rejection temperature. The \emph{characteristic equation} method answers
exactly that with three external temperatures and two regression constants.

Its basis is the D\"uhring diagram: for a given pair, the boiling temperature of the solution is very
nearly linear in the refrigerant saturation temperature at the same pressure. Following that line
through generator, condenser, absorber and evaporator collapses the cycle into a single driving
potential, the \emph{characteristic temperature difference}
\begin{equation}
\boxed{\;\Delta\Delta T = T_{\mathrm{gen}} - A\,T_{\mathrm{ac}} + (A-1)\,T_{\mathrm{evap}}\;}
\label{eq:spec:ddt}
\end{equation}
with $T_{\mathrm{gen}}$ the generator driving temperature, $T_{\mathrm{ac}}$ the common
condenser/absorber (rejection) temperature, $T_{\mathrm{evap}}$ the evaporator temperature and $A$ the
D\"uhring slope, $A\approx1.15$ for single-effect LiBr--water --- a \emph{typical literature value},
not a code constant. Because $A>1$, a kelvin lost at the rejection side costs more than a kelvin
gained at the generator side; Section~\ref{sec:spec:tower} turns that into an architectural argument.
The evaporator duty is linear in the potential:
\begin{equation}
\boxed{\;Q_{\mathrm{evap}} = a\,\Delta\Delta T + b\;}\qquad (a>0,\;b<0) .
\label{eq:spec:qevap}
\end{equation}
The constants $a$ and $b$ belong to one machine, obtained by regressing manufacturer part-load data,
and they absorb every internal approach the model does not resolve. Their signs carry meaning: $a$ is
the machine's gain, and the negative intercept $b$ encodes the internal losses that must be overcome
before any refrigerant is boiled, so \eqref{eq:spec:qevap} already contains its own shut-off,
\begin{equation}
\Delta\Delta T_{\min} = -\,b/a ,
\qquad Q_{\mathrm{evap}}\equiv 0 \ \text{ for } \Delta\Delta T \le \Delta\Delta T_{\min} .
\label{eq:spec:ddtmin}
\end{equation}
For a nominal \SI{100}{\kilo\watt} single-effect machine, $a\approx\SI{3.5}{\kilo\watt\per\kelvin}$
and $b\approx\SI{-30}{\kilo\watt}$ reproduce the nameplate at $\Delta\Delta T\approx\SI{37}{\kelvin}$.
\emph{Illustrative literature-scale figures for the reader's arithmetic --- not code constants and not
a product selection.}

\subsubsection{Generator side and the water-temperature drop}
Define the thermal COP and take the generator duty from it:
\begin{equation}
\mathrm{COP} = \frac{Q_{\mathrm{evap}}}{Q_{\mathrm{gen}}} ,
\qquad
Q_{\mathrm{gen}} = \frac{Q_{\mathrm{evap}}}{\mathrm{COP}} .
\label{eq:spec:qgen}
\end{equation}
Single-effect machines run $\mathrm{COP}=\numrange{0.65}{0.75}$; $0.70$ is the usual placeholder and
the value behind the resource tables of the topology report. \emph{Literature value.} The geothermal
stream feeding the generator then drops by
\begin{equation}
\Delta T_{\mathrm{gen,w}} = \frac{Q_{\mathrm{gen}}}{\md_{\mathrm{gen}}\,c_{p,\mathrm{w}}} ,
\qquad
T_{\mathrm{w,gen,out}} = T_{\mathrm{w,gen,in}} - \Delta T_{\mathrm{gen,w}} ,
\label{eq:spec:dtgen}
\end{equation}
and $T_{\mathrm{w,gen,out}}$ returns to the hydraulic mixing node exactly as the coil's
$T_{\mathrm{wh}}$ does. Design drops are \SIrange{8}{10}{\kelvin} with a maximum near
\SI{18}{\kelvin} --- \emph{literature values}, and the ones behind the flow-squeeze table of the
topology report.

One consistency point not to be skipped: $T_{\mathrm{gen}}$ in \eqref{eq:spec:ddt} is an
\emph{external} temperature, because $a$ and $b$ were regressed against external temperatures. Use
the mean of the generator water inlet and outlet,
\begin{equation}
T_{\mathrm{gen}} = \tfrac{1}{2}\left(T_{\mathrm{w,gen,in}} + T_{\mathrm{w,gen,out}}\right) ,
\label{eq:spec:tgenmean}
\end{equation}
and do \emph{not} add a separate internal approach --- that approach already lives in $b$. Since
\eqref{eq:spec:tgenmean} depends on $Q_{\mathrm{gen}}$, which depends on $Q_{\mathrm{evap}}$, which
depends on \eqref{eq:spec:ddt}, the generator side is itself a small fixed point converging in two or
three passes.

\subsubsection{Heat rejected}
Steady-state balance on the machine, solution-pump work neglected (well under \SI{1}{\percent} of
$Q_{\mathrm{gen}}$):
\begin{equation}
Q_{\mathrm{rej}} = Q_{\mathrm{gen}} + Q_{\mathrm{evap}}
 = Q_{\mathrm{evap}}\left(1+\frac{1}{\mathrm{COP}}\right)
 \approx 2.4\,Q_{\mathrm{evap}} \ \text{ at } \mathrm{COP}=0.70 .
\label{eq:spec:qrej}
\end{equation}
This is the number that sizes the tower, and it is worth pausing on: an absorption machine rejects
roughly \emph{two and a half times} its useful cooling, against about $1.25$ for a
vapour-compression chiller. The tower is correspondingly large, and so is its water consumption. That
is the price of being driven by heat the well supplies for nothing.

\subsubsection{The two limits, and why this component binds the architecture}
Equation~\eqref{eq:spec:qevap} is smooth and will return a duty at any temperature handed to it. The
physical machine will not. Two bounds must be imposed \emph{outside} the characteristic equation,
because neither is contained in it:
\begin{description}[leftmargin=1.6em,nosep]
\item[Lower bound, $T_{\mathrm{gen}}\gtrsim\SI{78}{\celsius}$.] Below this the generator cannot raise
the solution above its boiling point at condenser pressure and no refrigerant vapour is produced. The
machine does not degrade gracefully; it stops. Implement as a hard gate returning
$Q_{\mathrm{evap}}=0$ with an explicit warning, in the spirit of the wheel's regime diagnostic
(Section~\ref{sec:dw:regime}) --- a silent zero is worse than none.
\item[Upper bound, $T_{\mathrm{gen}}\lesssim\SI{110}{\celsius}$.] Above this the concentration
required at the generator exit approaches the crystallisation line, and the corrosion rate of carbon
steel in hot concentrated LiBr rises steeply. The penalty is not lost performance but a blocked
machine and a maintenance event. Implement as a clamp with a warning, never as an extrapolation.
\end{description}
Both thresholds are \emph{typical literature values} for hot-water-fired single-effect machines;
different manufacturers quote \SIrange{75}{80}{\celsius} and \SIrange{110}{120}{\celsius}. They are
not code constants, and the implementation should carry them as named configuration parameters so a
real data sheet can replace them in one place.

\paragraph{The architectural consequence, stated plainly.} Of every component in the master
architecture, \emph{this one alone requires the high-temperature tier}. The desiccant wheel
regenerates at roughly \SI{66}{\celsius}; the winter heating duty needs only the supply-air
temperature; the humidifier and the IEC need no geothermal water at all. Only the absorption
generator asks for \SI{78}{\celsius} or better, and only modes M3 and M5 contain it. As the reservoir
depletes and the return temperature drifts down, the plant therefore does not lose capability
uniformly: it crosses the \SI{78}{\celsius} gate first, and at that moment M3 and M5 switch off while
every other mode continues untouched. \emph{The serviceable region of the load plane shrinks from the
cooling corner first.} That is the strongest system-level result the architecture yields, and it is
entirely a consequence of the lower bound above --- which is why the bound must be a declared
parameter and not a magic number buried in a branch.

%---------------------------------------------------------------------
\subsection{Wet cooling tower (serves the chiller)}
\label{sec:spec:tower}

\emph{Status: specification only. Not implemented in the current model version.}

\subsubsection{Hypothesis: Merkel, and why the potential is enthalpy}
Water falls through fill against an upward air stream and part of it evaporates. Heat and mass
transfer occur together, so a temperature difference is not the correct driving potential --- a tower
can cool water \emph{below} the ambient dry bulb, which no temperature-driven exchanger can do.
Merkel's simplification resolves this: assume the interface air is saturated at the local bulk water
temperature and that the Lewis number is unity, and the coupled problem collapses to one transfer
equation driven by an \emph{enthalpy} difference. The tower then admits the same
$\varepsilon$--$\mathrm{NTU}$ treatment as the regeneration coil, with enthalpy in place of
temperature.

\subsubsection{Equations}
Define the saturated-air enthalpy function, built from the notebook's own $\psat$ and enthalpy
correlations,
\begin{equation}
h_{s}(T) \equiv h\!\left(T,w_{s}(T)\right)
= 1.006\,T + \frac{0.622\,\psat(T)}{\patm-\psat(T)}\left(2501+1.86\,T\right) ,
\label{eq:spec:hsat}
\end{equation}
and its secant slope over the operating range, which plays the role of a specific heat on the water
side:
\begin{equation}
C_{s} = \frac{h_{s}(T_{\mathrm{t,i}}) - h_{s}(T_{\mathrm{t,o}})}
             {T_{\mathrm{t,i}} - T_{\mathrm{t,o}}} .
\label{eq:spec:Cs}
\end{equation}
The water stream's \emph{equivalent} capacity rate --- its capacity expressed in enthalpy units so it
can be compared with the air's --- is $\md_{\mathrm{t}}c_{p,\mathrm{w}}/C_{s}$, while the air's is
simply $\md_{\mathrm{a,t}}$, the potential already being an enthalpy. Hence
\begin{equation}
C_{r} = \frac{\md_{\mathrm{a,t}}\,C_{s}}{\md_{\mathrm{t}}\,c_{p,\mathrm{w}}} ,
\qquad
\mathrm{NTU}_{\mathrm{t}} = \frac{\UA_{\mathrm{t}}}{\md_{\mathrm{a,t}}} ,
\label{eq:spec:towerNTU}
\end{equation}
with $\UA_{\mathrm{t}}$ a mass-transfer conductance carrying the units of a mass flow. The air-side
effectiveness uses the counterflow relation already coded for the coil,
\begin{equation}
\varepsilon_{\mathrm{t}} =
\frac{1-\exp\!\left[-\mathrm{NTU}_{\mathrm{t}}(1-C_{r})\right]}
     {1-C_{r}\exp\!\left[-\mathrm{NTU}_{\mathrm{t}}(1-C_{r})\right]} ,
\qquad
\varepsilon_{\mathrm{t}} = \frac{\mathrm{NTU}_{\mathrm{t}}}{1+\mathrm{NTU}_{\mathrm{t}}}
\ \text{ if } C_{r}=1 ,
\label{eq:spec:towereps}
\end{equation}
and the duty is the effectiveness times the maximum enthalpy potential --- the gap between saturation
at the hot water inlet and the actual inlet air:
\begin{equation}
Q_{\mathrm{t}} = \varepsilon_{\mathrm{t}}\,\md_{\mathrm{a,t}}
\left[h_{s}(T_{\mathrm{t,i}}) - h_{\mathrm{a,i}}\right] .
\label{eq:spec:towerQ}
\end{equation}
Since $h_{\mathrm{a,i}} = h_{s}(T_{\mathrm{a,i}}^{\mathrm{wb}})$ to within the makeup term, the
potential in \eqref{eq:spec:towerQ} is measured from the \emph{ambient wet bulb} --- the analytical
statement of the fact that a tower chases the wet bulb and not the dry bulb. Outlet temperature,
range and approach follow:
\begin{equation}
T_{\mathrm{t,o}} = T_{\mathrm{t,i}} - \frac{Q_{\mathrm{t}}}{\md_{\mathrm{t}}c_{p,\mathrm{w}}} ,
\qquad
\mathrm{range} = T_{\mathrm{t,i}} - T_{\mathrm{t,o}} ,
\qquad
\mathrm{approach} = T_{\mathrm{t,o}} - T_{\mathrm{a,i}}^{\mathrm{wb}} > 0 .
\label{eq:spec:towerT}
\end{equation}
Equations \eqref{eq:spec:Cs}--\eqref{eq:spec:towerT} are implicit in $T_{\mathrm{t,o}}$ through
$C_{s}$, so the block iterates internally; two or three passes from
$T_{\mathrm{t,o}}=T_{\mathrm{a,i}}^{\mathrm{wb}}+\SI{4}{\kelvin}$ suffice.

Water leaves the tower three ways --- evaporation, from the air-side moisture pickup with the outlet
air taken as saturated at its own enthalpy (the usual Merkel closure); blowdown, which removes
dissolved solids; and drift, which is mechanical carryover:
\begin{equation}
\md_{\mathrm{evap}} = \md_{\mathrm{a,t}}\left(w_{\mathrm{a,o}}-w_{\mathrm{a,i}}\right) ,
\qquad
\md_{\mathrm{blow}} = \frac{\md_{\mathrm{evap}}}{N_{c}-1} ,
\qquad
\md_{\mathrm{makeup}} = \md_{\mathrm{evap}} + \md_{\mathrm{blow}} + \md_{\mathrm{drift}} .
\label{eq:spec:towerwater}
\end{equation}
A useful implementation check: $\md_{\mathrm{evap}}$ should come out near \SI{1}{\percent} of the
circulating flow for every \SI{6}{\kelvin} of range, which follows from $c_{p,\mathrm{w}}/\hfg$ and
the fact that roughly three-quarters of a wet tower's duty is latent.

\paragraph{Typical values.} Design approach \SIrange{3}{5}{\kelvin}; design range
\SIrange{5}{8}{\kelvin}; $\mathrm{NTU}_{\mathrm{t}}=\numrange{1.5}{2.5}$; cycles of concentration
$N_{c}=\numrange{3}{5}$; drift \SIrange{0.005}{0.02}{\percent} of circulating flow. \emph{All
literature ranges for mechanical-draught counterflow towers; none is a code constant.}

\subsubsection{Why the tower cannot be solved after the chiller}
Here is the statement this subsection exists to make. The tower's approach sets $T_{\mathrm{t,o}}$;
the chiller rejects into that water, so its condenser/absorber temperature is
\begin{equation}
T_{\mathrm{ac}} = T_{\mathrm{t,o}} + \delta_{\mathrm{ac}} ,
\label{eq:spec:Tac}
\end{equation}
with $\delta_{\mathrm{ac}}$ the rejection-side approach inside the machine
(\SIrange{3}{6}{\kelvin}, \emph{literature}). But $T_{\mathrm{ac}}$ is exactly the term the D\"uhring
slope multiplies in \eqref{eq:spec:ddt}, so
\begin{equation}
\frac{\partial Q_{\mathrm{evap}}}{\partial T_{\mathrm{t,o}}} = -\,a\,A \approx
\SI{-4}{\kilo\watt\per\kelvin}
\qquad\text{against}\qquad
\frac{\partial Q_{\mathrm{evap}}}{\partial T_{\mathrm{gen}}} = a \approx
\SI{3.5}{\kilo\watt\per\kelvin}
\label{eq:spec:sens}
\end{equation}
at the illustrative constants above. \emph{A kelvin surrendered at the cooling tower costs more
evaporator duty than a kelvin gained at the well returns.} And the loop closes: $Q_{\mathrm{evap}}$
sets $Q_{\mathrm{rej}}$ by \eqref{eq:spec:qrej}, which sets the tower duty, which sets
$T_{\mathrm{t,o}}$, which sets $T_{\mathrm{ac}}$, which sets $Q_{\mathrm{evap}}$. Solving the tower
after the chiller --- fixing $T_{\mathrm{ac}}$ at a nominal number and sizing the tower to whatever
falls out --- is not a mild approximation. It discards the dominant feedback and reports capacity the
plant does not have, on exactly the humid design days when the wet bulb is high and the error is
largest.

The two must therefore be closed simultaneously. The cheapest correct implementation is a
one-dimensional residual in $T_{\mathrm{t,o}}$,
\begin{equation}
R(T_{\mathrm{t,o}}) = T_{\mathrm{t,o}}
 - \left[T_{\mathrm{t,i}} - \frac{Q_{\mathrm{t}}}{\md_{\mathrm{t}}c_{p,\mathrm{w}}}\right] = 0 ,
\qquad
T_{\mathrm{t,i}} = T_{\mathrm{t,o}} + \frac{Q_{\mathrm{rej}}(T_{\mathrm{t,o}})}
{\md_{\mathrm{t}}c_{p,\mathrm{w}}} ,
\label{eq:spec:coupled}
\end{equation}
driven by bisection on the physically bounded interval
$T_{\mathrm{a,i}}^{\mathrm{wb}} < T_{\mathrm{t,o}} < T_{\mathrm{t,i}}$. It is monotone, it cannot
leave the interval, and it costs perhaps a dozen evaluations of two algebraic blocks --- negligible
beside the DBHE march.

%---------------------------------------------------------------------
\subsection{Implementation order, and one unresolved defect}
\label{sec:spec:order}

\emph{Status: none of the following is implemented.}

Table~\ref{tab:spec:order} gives the sequence, ordered by unlocked capability per unit of effort. It
is not arbitrary: item~0 is a prerequisite for 1, 2 and 4, and items 3--5 cannot be separated,
because Section~\ref{sec:spec:tower} showed that chiller and tower share a fixed point.

\begin{table}[htbp]
\centering
\caption{Implementation order for the pending components. ``Unlocks'' refers to the modes of the
master architecture. Nothing in this table is implemented.}
\label{tab:spec:order}
\begin{tabular}{@{}clll@{}}
\toprule
\# & Item & Unlocks & Note \\
\midrule
0 & Wet-bulb and dew-point solvers & --- & Eq.~\eqref{eq:spec:wbdp}; not a component \\
1 & Adiabatic humidifier & M4 (and M8) & One block; needs no geothermal water \\
2 & Indirect evaporative cooler & M2 & One block plus a saturated-outlet solve \\
3 & Chilled-water loop & --- & Hydraulic; prerequisite for 4 \\
4 & Absorption chiller & M3, M5 & Must be solved with 5 \\
5 & Wet cooling tower & M3, M5 & Must be solved with 4 \\
\bottomrule
\end{tabular}
\end{table}

Item~1 is the highest-value single change in the backlog: a dozen lines, no new loop, and it opens
the one operating mode that is independent of the reservoir and therefore immune to depletion. Items
4 and 5 are the largest change, and they open the two modes that are the \emph{first} to be lost to
depletion. That asymmetry is worth keeping in view when the work is scheduled.

\subsubsection{The unresolved defect: process ordering in M5}
Mode M5 activates the absorption chiller with a dry chilled-water coil \emph{and} the humidifier, to
serve deep cooling-plus-humidification loads that M4 cannot reach. In the architecture as drawn, the
humidifier sits on prong~E, \emph{upstream} of the chilled-water coil in fork~2. The air is therefore
humidified before it is cooled. This is the wrong order, for three reasons:
\begin{enumerate}[nosep]
\item \textbf{It may saturate.} Adiabatic saturation drives the air toward $\RH=1$ along a
constant-enthalpy line. Presenting a nearly saturated stream to a cold coil face guarantees
condensation on the fins.
\item \textbf{It undoes itself.} Whatever condenses on the coil is moisture the humidifier has just
evaporated. The plant pays evaporative duty to add water and chiller duty to take the same water back
out --- latent work performed twice for a net result of zero.
\item \textbf{It wastes the free cooling.} The humidifier's own sensible cooling,
Eq.~\eqref{eq:spec:humduty}, is delivered upstream of the coil, where the coil simply sees a colder
inlet and throttles back. Downstream of the coil the same cooling would \emph{add} to delivered
capacity instead of substituting for it.
\end{enumerate}

\emph{Cool first, then humidify} is the correct sequence: the coil brings the air down at nearly
constant $w$, and the humidifier then walks it up the constant-wet-bulb line to the supply state,
its own evaporative cooling counting toward the sensible duty rather than against it.
\textbf{The architecture does not currently provide that path} --- prong~E has no branch that
reaches the humidifier after fork~2.

Three remedies exist, and the choice is a genuine open question rather than an oversight to be
quietly patched:
\begin{itemize}[nosep]
\item \emph{Reposition.} Add a duct branch placing the humidifier downstream of the chilled-water
coil. Physically correct, and the only option that captures the humidifier's cooling as useful duty.
Costs a new prong and one more damper in a valve network that is already the most intricate part of
the plant.
\item \emph{Duplicate.} Install humidification capacity in both positions and select by mode. Correct
in every mode, but it buys generality with hardware idle most of the year.
\item \emph{Detune.} Keep the present order and cap $\varepsilon_{\mathrm{sat}}$ so the coil face
never sees $\RH>\RH_{\max}$. A pure control workaround: it removes the condensation risk but keeps
defects 2 and 3, and it forfeits latent capacity precisely in the loads that forced M5 in the first
place.
\end{itemize}
Whichever is adopted, the humidifier block of Section~\ref{sec:spec:hum} is unchanged --- only its
position in the stream sequence moves. That is the advantage of specifying components as pure
functions of their inlet states, and it is why this defect is a topology question and not a modelling
one.
